\documentclass[10pt,a4paper]{amsart}
\usepackage{amsmath,amssymb,amsthm,graphicx}
\usepackage{verbatim}
\newenvironment{reference}{\comment}{\endcomment}
%\newenvironment{reference}{}{}
\newenvironment{slogan}{\comment}{\endcomment}
\newenvironment{history}{\comment}{\endcomment}

% For commutative diagrams we use Xy-pic
\usepackage[all]{xy}

% We use 2cell for 2-commutative diagrams.
\xyoption{2cell}
% Theorem environments.
%
\theoremstyle{plain}
\newtheorem{theorem}[subsection]{Theorem}
\newtheorem{proposition}[subsection]{Proposition}
\newtheorem{lemma}[subsection]{Lemma}

\theoremstyle{definition}
\newtheorem{definition}[subsection]{Definition}
\newtheorem{example}[subsection]{Example}
\newtheorem{exercise}[subsection]{Exercise}
\newtheorem{situation}[subsection]{Situation}

\theoremstyle{remark}
\newtheorem{remark}[subsection]{Remark}
\newtheorem{remarks}[subsection]{Remarks}

\numberwithin{equation}{subsection}

% Macros
%
\def\lim{\mathop{\mathrm{lim}}\nolimits}
\def\colim{\mathop{\mathrm{colim}}\nolimits}
\def\Spec{\mathop{\mathrm{Spec}}}
\def\Hom{\mathop{\mathrm{Hom}}\nolimits}
\def\Ext{\mathop{\mathrm{Ext}}\nolimits}
\def\SheafHom{\mathop{\mathcal{H}\!\mathit{om}}\nolimits}
\def\SheafExt{\mathop{\mathcal{E}\!\mathit{xt}}\nolimits}
\def\Sch{\mathit{Sch}}
\def\Mor{\mathop{\mathrm{Mor}}\nolimits}
\def\Ob{\mathop{\mathrm{Ob}}\nolimits}
\def\Sh{\mathop{\mathit{Sh}}\nolimits}
\def\NL{\mathop{N\!L}\nolimits}
\def\CH{\mathop{\mathrm{CH}}\nolimits}
\def\proetale{{pro\text{-}\acute{e}tale}}
\def\etale{{\acute{e}tale}}
\def\QCoh{\mathit{QCoh}}
\def\Ker{\mathop{\mathrm{Ker}}}
\def\Im{\mathop{\mathrm{Im}}}
\def\Coker{\mathop{\mathrm{Coker}}}
\def\Coim{\mathop{\mathrm{Coim}}}

% Boxtimes
%
\DeclareMathSymbol{\boxtimes}{\mathbin}{AMSa}{"02}

%
% Macros for moduli stacks/spaces
%
\def\QCohstack{\mathcal{QC}\!\mathit{oh}}
\def\Cohstack{\mathcal{C}\!\mathit{oh}}
\def\Spacesstack{\mathcal{S}\!\mathit{paces}}
\def\Quotfunctor{\mathrm{Quot}}
\def\Hilbfunctor{\mathrm{Hilb}}
\def\Curvesstack{\mathcal{C}\!\mathit{urves}}
\def\Polarizedstack{\mathcal{P}\!\mathit{olarized}}
\def\Complexesstack{\mathcal{C}\!\mathit{omplexes}}
% \Pic is the operator that assigns to X its picard group, usage \Pic(X)
% \Picardstack_{X/B} denotes the Picard stack of X over B
% \Picardfunctor_{X/B} denotes the Picard functor of X over B
\def\Pic{\mathop{\mathrm{Pic}}\nolimits}
\def\Picardstack{\mathcal{P}\!\mathit{ic}}
\def\Picardfunctor{\mathrm{Pic}}
\def\Deformationcategory{\mathcal{D}\!\mathit{ef}}
\usepackage{cleveref}
\crefname{originallemma}{original lemma}{original lemmas}
\Crefname{originallemma}{Original lemma}{Original lemmas}
\makeatletter
\newlabel{lemma-abelian-variety-projective}{{0BFA}{source}{}{sourcepointer-1}{}}
\newlabel{lemma-abelian-variety-projective@cref}{{[originallemma][0][]0BFA}{[1][0][]source}{}{}{}}
\makeatother
\crefname{originallemma}{original lemma}{original lemmas}
\Crefname{originallemma}{Original lemma}{Original lemmas}
\makeatletter
\newlabel{lemma-abelian-variety-abelian}{{0BFD}{source}{}{sourcepointer-2}{}}
\newlabel{lemma-abelian-variety-abelian@cref}{{[originallemma][0][]0BFD}{[1][0][]source}{}{}{}}
\makeatother
\crefname{originallemma}{original lemma}{original lemmas}
\Crefname{originallemma}{Original lemma}{Original lemmas}
\makeatletter
\newlabel{lemma-degree-multiplication-by-d}{{0BFG}{source}{}{sourcepointer-3}{}}
\newlabel{lemma-degree-multiplication-by-d@cref}{{[originallemma][0][]0BFG}{[1][0][]source}{}{}{}}
\makeatother
\crefname{originallemma}{original lemma}{original lemmas}
\Crefname{originallemma}{Original lemma}{Original lemmas}
\makeatletter
\newlabel{lemma-abelian-variety-multiplication-by-d-etale}{{0BFH}{source}{}{sourcepointer-4}{}}
\newlabel{lemma-abelian-variety-multiplication-by-d-etale@cref}{{[originallemma][0][]0BFH}{[1][0][]source}{}{}{}}
\makeatother
\crefname{originalreference}{original reference}{original references}
\Crefname{originalreference}{Original reference}{Original references}
\makeatletter
\newlabel{morphisms-lemma-flat-unramified-etale}{{02GV}{source}{}{sourcepointer-5}{}}
\newlabel{morphisms-lemma-flat-unramified-etale@cref}{{[originalreference][0][]02GV}{[1][0][]source}{}{}{}}
\makeatother
\crefname{originalreference}{original reference}{original references}
\Crefname{originalreference}{Original reference}{Original references}
\makeatletter
\newlabel{morphisms-lemma-set-points-where-fibres-unramified}{{0475}{source}{}{sourcepointer-6}{}}
\newlabel{morphisms-lemma-set-points-where-fibres-unramified@cref}{{[originalreference][0][]0475}{[1][0][]source}{}{}{}}
\makeatother
\crefname{originallemma}{original lemma}{original lemmas}
\Crefname{originallemma}{Original lemma}{Original lemmas}
\makeatletter
\newlabel{lemma-abelian-variety-multiplication-by-p}{{0C0Y}{source}{}{sourcepointer-7}{}}
\newlabel{lemma-abelian-variety-multiplication-by-p@cref}{{[originallemma][0][]0C0Y}{[1][0][]source}{}{}{}}
\makeatother
\title{Review of eight standard properties of an abelian variety}
\author{The Stacks Project Authors}
\date{stacks-project:groupoids.tex:local-snapshot-20261012; GNU FDL 1.2 or later}
\begin{document}
\maketitle
\noindent\textit{Editorial note.} All delivered mathematical text is the author's own, copied verbatim from the pinned Stacks Project chapter file groupoids.tex; no statement or proof step was rewritten, reordered or removed.\par\medskip
\noindent\textit{Editorial note.} Span recomputation by label pairing (LANE\_BRIEF 10.2): in this chapter the author's proof environment follows the statement environment as a standalone block; the spec covers the complete environment pair (statement span, proof span) as recorded in the pool candidate.\par\medskip
\noindent\textit{Editorial note.} Occupation scan: this chapter's pool has 47 of 128 labels already delivered by the S7 line (S7-gp-* registrations, statement\_anchor-prose scan; withdrawn cards excluded); none of the 11 accepted labels of this batch is among them (gp-occupied.json, 2026-10-14).\par\medskip
\ifdefined\hypertarget\hypertarget{sourcepointer-1}{}\fi
\noindent\textit{Source reference.} Original lemma 0BFA: sources/stacks-chow/tags.txt, source lines 14378--14378. Cross-chapter standard prerequisite cited by the author inside the counted span; it is used as a pointer and its own proof is not part of this item.\par\medskip
\ifdefined\hypertarget\hypertarget{sourcepointer-2}{}\fi
\noindent\textit{Source reference.} Original lemma 0BFD: sources/stacks-chow/tags.txt, source lines 14381--14381. Cross-chapter standard prerequisite cited by the author inside the counted span; it is used as a pointer and its own proof is not part of this item.\par\medskip
\ifdefined\hypertarget\hypertarget{sourcepointer-3}{}\fi
\noindent\textit{Source reference.} Original lemma 0BFG: sources/stacks-chow/tags.txt, source lines 14384--14384. Cross-chapter standard prerequisite cited by the author inside the counted span; it is used as a pointer and its own proof is not part of this item.\par\medskip
\ifdefined\hypertarget\hypertarget{sourcepointer-4}{}\fi
\noindent\textit{Source reference.} Original lemma 0BFH: sources/stacks-chow/tags.txt, source lines 14385--14385. Cross-chapter standard prerequisite cited by the author inside the counted span; it is used as a pointer and its own proof is not part of this item.\par\medskip
\ifdefined\hypertarget\hypertarget{sourcepointer-5}{}\fi
\noindent\textit{Source reference.} Original reference 02GV: sources/stacks-chow/tags.txt, source lines 3242--3242. Cross-chapter standard prerequisite cited by the author inside the counted span; it is used as a pointer and its own proof is not part of this item.\par\medskip
\ifdefined\hypertarget\hypertarget{sourcepointer-6}{}\fi
\noindent\textit{Source reference.} Original reference 0475: sources/stacks-chow/tags.txt, source lines 5496--5496. Cross-chapter standard prerequisite cited by the author inside the counted span; it is used as a pointer and its own proof is not part of this item.\par\medskip
\ifdefined\hypertarget\hypertarget{sourcepointer-7}{}\fi
\noindent\textit{Source reference.} Original lemma 0C0Y: sources/stacks-chow/tags.txt, source lines 15103--15103. Cross-chapter standard prerequisite cited by the author inside the counted span; it is used as a pointer and its own proof is not part of this item.\par\medskip
\begin{proposition}
\label{proposition-review-abelian-varieties}
\begin{reference}
Wonderfully explained in \cite{AVar}.
\end{reference}
Let $A$ be an abelian variety over a field $k$. Then
\begin{enumerate}
\item $A$ is projective over $k$,
\item $A$ is a commutative group scheme,
\item the morphism $[n] : A \to A$ is surjective for all $n \geq 1$,
\item if $k$ is algebraically closed, then $A(k)$ is a divisible abelian group,
\item $A[n] = \Ker([n] : A \to A)$ is a finite group scheme of degree
$n^{2\dim A}$ over $k$,
\item $A[n]$ is \'etale over $k$ if and only if $n \in k^*$,
\item if $n \in k^*$ and $k$ is algebraically closed,
then $A(k)[n] \cong (\mathbf{Z}/n\mathbf{Z})^{\oplus 2\dim(A)}$,
\item if $k$ is algebraically closed of characteristic $p > 0$, then
there exists an integer $0 \leq f \leq \dim(A)$ such that
$A(k)[p^m] \cong (\mathbf{Z}/p^m\mathbf{Z})^{\oplus f}$
for all $m \geq 1$.
\end{enumerate}
\end{proposition}

\begin{proof}
Part (1) follows from Lemma \ref{lemma-abelian-variety-projective}.
Part (2) follows from Lemma \ref{lemma-abelian-variety-abelian}.
Part (3) follows from Lemma \ref{lemma-degree-multiplication-by-d}.
If $k$ is algebraically closed then surjective morphisms of varieties
over $k$ induce surjective maps on $k$-rational points, hence
(4) follows from (3).
Part (5) follows from Lemma \ref{lemma-degree-multiplication-by-d}
and the fact that a base change of a finite locally free morphism
of degree $N$ is a finite locally free morphism of degree $N$.
Part (6) follows from 
Lemma \ref{lemma-abelian-variety-multiplication-by-d-etale}.
Namely, if $n$ is invertible in $k$, then $[n]$ is \'etale
and hence $A[n]$ is \'etale over $k$.
On the other hand, if $n$ is not invertible in $k$, then
$[n]$ is not \'etale at $e$ and it follows that $A[n]$
is not \'etale over $k$ at $e$ (use
Morphisms, Lemmas \ref{morphisms-lemma-flat-unramified-etale} and
\ref{morphisms-lemma-set-points-where-fibres-unramified}).

\medskip\noindent
Assume $k$ is algebraically closed. Set $g = \dim(A)$. Proof of (7).
Let $\ell$ be a prime number which is invertible in $k$. Then we see that
$$
A[\ell](k) = A(k)[\ell]
$$
is a finite abelian group, annihilated by $\ell$, of order $\ell^{2g}$.
It follows that it is isomorphic to $(\mathbf{Z}/\ell\mathbf{Z})^{2g}$
by the structure theory for finite abelian groups. Next, we consider
the short exact sequence
$$
0 \to A(k)[\ell] \to A(k)[\ell^2] \xrightarrow{\ell} A(k)[\ell] \to 0
$$
Arguing similarly as above we conclude that 
$A(k)[\ell^2] \cong (\mathbf{Z}/\ell^2\mathbf{Z})^{2g}$.
By induction on the exponent we find that
$A(k)[\ell^m] \cong (\mathbf{Z}/\ell^m\mathbf{Z})^{2g}$.
For composite integers $n$ prime to the characteristic of $k$
we take primary parts and we find the correct shape of the
$n$-torsion in $A(k)$.
The proof of (8) proceeds in exactly the same way, using that
Lemma \ref{lemma-abelian-variety-multiplication-by-p} gives
$A(k)[p] \cong (\mathbf{Z}/p\mathbf{Z})^{\oplus f}$ for some $0 \leq f \leq g$.
\end{proof}

\end{document}
